% 临时打印页：只把板子里的某一小块单独排成两三页的 PDF，方便单模块补印，
% 不用为了改一个类就把近 200 页的 template-main.pdf 整本重打。
%
% 要印的片段在同目录 fragment.tex（两个 subsection），./build.sh 编译，产物 out/vandermonde_temporary_print.pdf。
% 一个临时打印页一个文件夹，规范见 TemporaryPrint/README.md。
% preamble 是从 template-main.tex 裁出来的最小可用子集（只留这类代码页真正用到的宏包），
% 字体 / 页边距 / minted 设置与正式板子一致，但**排成单栏**：
% 补印页只是拿来看一个模块，单栏行更长、代码不用挤，比照搬书里的双栏好读。

\documentclass{article}

\usepackage{amsmath}
\usepackage{amssymb}
\usepackage{csquotes}

\usepackage{xeCJK}
\setCJKmainfont[BoldFont={Noto Serif CJK SC Bold}, ItalicFont={Kaiti SC}]{Noto Serif CJK SC}
\setCJKsansfont[BoldFont={Noto Sans CJK SC Bold}]{Noto Sans CJK SC}
\setCJKmonofont{Noto Sans CJK SC}

\usepackage[table]{xcolor}
% cachedir 与主文档分开，避免两边的 minted 缓存互相打架
\usepackage[cachedir=_minted-temporary-print]{minted}
% 片段里有 TikZ 示意图（中位数例题），与主文档同样加载 tikz + brace 装饰
\usepackage{tikz}
\usetikzlibrary{decorations.pathreplacing}
\usepackage[a4paper, left=2cm, right=2cm, top=2.5cm, bottom=2.5cm]{geometry}

% 与 template-main.tex 逐字一致的 minted 设置
\setminted{
	breaklines=true,
	breakanywhere=true,
	fontsize=\small,
	frame=single,
	bgcolor=gray!5,
	tabsize=2,
}
\setmintedinline{breakanywhere=false}

\usepackage[colorlinks=true, linkcolor=black, anchorcolor=black, citecolor=black, filecolor=black, menucolor=black, runcolor=black, urlcolor=blue]{hyperref}

% 单块补印不要 "0.1" 这种脱离全书的孤零零编号
\setcounter{secnumdepth}{-1}

\begin{document}

	\begin{center}
		{\Large \textbf{临时打印页：计数与组合 · 范德蒙德卷积（含中位数例题）与常用组合恒等式}} \\[4pt]
		{\small 补印自 \texttt{sections/04\_count\_combinatorics.tex}\quad·\quad 2026-10-08}
	\end{center}

	\vspace{0.5em}

	% 补印自 sections/04_count_combinatorics.tex 的「范德蒙德卷积」（含中位数为 X 例题）与「常用组合恒等式」两个 subsection（2026-10-08）
% 正式板子是唯一真身：改动先改 sections/，再复制到这里

		\subsection{范德蒙德卷积}

		\[
			\sum_{i}\binom{m}{i}\binom{n}{k-i} = \binom{m+n}{k}.
		\]
		组合意义：从 $m$ 男 $n$ 女里选 $k$ 人，按选了几个男的分类。两个常用变体：
		\[
			\sum_{i}\binom{m}{i}\binom{n}{i} = \binom{m+n}{m},
			\qquad
			\sum_{i}\binom{n}{i}^{2} = \binom{2n}{n}.
		\]
		（第一式把 $\binom{n}{i}$ 换成 $\binom{n}{n-i}$ 即可。）
		见到「两个组合数相乘再对下标求和」就应该先试着往这里凑。

		\textbf{「左右各选 $r$ 个」式}：左边 $a$ 个、右边 $b$ 个，两边选\textbf{一样多}的方案总数
		\[
			\sum_{r}\binom{a}{r}\binom{b}{r} = \binom{a+b}{a}.
		\]
		一一对应：「左边选中的」并上「右边\textbf{没}选中的」恰好是 $a$ 个，反过来也唯一还原。
		典型场景：数组排序后\textbf{固定中位数所在位置}，两侧必须等量选。
		两侧数量差恰好为 $d$ 时同理（平移版）：
		\[
			\sum_{r}\binom{a}{r}\binom{b}{r+d} = \binom{a+b}{a+d}.
		\]

		\subsubsection{例：中位数为 X（2026 CCPC 网络赛 M）}

		给长为 $n$ 的序列和 $x$，数中位数恰好为 $x$ 的非空子序列个数（偶数个取中间两数平均；不同下标集合算不同）。$n\le 2\times10^5$。

		\begin{center}
			\begin{tikzpicture}[x=0.62cm, y=0.62cm, font=\scriptsize]
				\tikzset{c/.style={draw, minimum size=0.62cm, inner sep=0pt},
					sel/.style={c, fill=blue!25}, mid/.style={c, fill=red!35},
					ban/.style={c, fill=gray!35}}
				% 奇数：中心 i = 3，左 3 个选 2 个、右 4 个选 2 个
				\node[left] at (-0.5, 0) {奇};
				\foreach \k/\s in {0/sel, 1/c, 2/sel, 3/mid, 4/c, 5/sel, 6/c, 7/sel}
					\node[\s] at (\k, 0) {\k};
				\draw[decorate, decoration={brace, mirror, amplitude=4pt}] (-0.45, -0.45) -- (2.45, -0.45)
					node[midway, below=4pt] {$i$ 个里选 $r$};
				\draw[decorate, decoration={brace, mirror, amplitude=4pt}] (3.55, -0.45) -- (7.45, -0.45)
					node[midway, below=4pt] {$n-1-i$ 个里选 $r$};
				% 偶数：中心对 i = 2, j = 5，中间不能选
				\node[left] at (-0.5, -2.4) {偶};
				\foreach \k/\s in {0/c, 1/sel, 2/mid, 3/ban, 4/ban, 5/mid, 6/sel, 7/c}
					\node[\s] at (\k, -2.4) {\k};
				\draw[decorate, decoration={brace, mirror, amplitude=4pt}] (-0.45, -2.85) -- (1.45, -2.85)
					node[midway, below=4pt] {$i$ 个里选 $r$};
				\draw[decorate, decoration={brace, mirror, amplitude=4pt}] (5.55, -2.85) -- (7.45, -2.85)
					node[midway, below=4pt] {$n-1-j$ 个里选 $r$};
			\end{tikzpicture}

			\vspace{2pt}
			{\footnotesize\sffamily\color{gray} 图注：排序后 $n=8$ 的下标。\textcolor{red!70!black}{红} = 中心，\textcolor{blue!70!black}{蓝} = 选中，灰 = 偶数情形两中心之间不能选。奇数中心 $i=3$ 两侧各选 $r=2$；偶数中心对 $(2,5)$ 两侧各选 $r=1$。}
		\end{center}

		\begin{enumerate}
			\item \textbf{去分数、恰好转不超过}：比 $2\cdot\mathrm{med}$（必为整数）。设 $F(Y)$ 为 $2\cdot\mathrm{med}\le Y$ 的集合数，答案 $=F(2x)-F(2x-1)$。
			\item \textbf{排序后按中心计数}（重复值也当不同位置，下标集合自然区分）：
			      奇数中心 $i$（$0$-based）：$\sum_r\binom{i}{r}\binom{n-1-i}{r}=\binom{n-1}{i}$，条件 $2a_i\le Y$；
			      偶数中心对 $i<j$：$\sum_r\binom{i}{r}\binom{n-1-j}{r}=\binom{i+n-1-j}{i}$，条件 $a_i+a_j\le Y$。
			\item \textbf{双指针 + 朱世杰}：$i$ 增大时最大合法 $j$（记 $q$）单调不增；一整段 $j\in(i,q]$ 合并为
			      \[
			      	\begin{aligned}
			      		&\sum_{j=i+1}^{q}\binom{i+n-1-j}{i}\\
			      		={}&\binom{n-1}{i+1}-\binom{i+n-q-1}{i+1}.
			      	\end{aligned}
			      \]
		\end{enumerate}
		总复杂度 $O(n\log n)$（排序）。

		\begin{minted}{cpp}
// 依赖上文 namespace Comb（CC），mod = 998244353，MAXN >= n
// F(Y) = 「2·中位数 <= Y」的非空下标集合数，a 已排序
ll F(const vector<ll> &a, ll Y) {
	int n = SZ(a); ll res = 0;
	for (int i = 0; i < n; i++)        // 奇数：中心 i，两侧各选 r 个
		if (2 * a[i] <= Y) res += CC(n - 1, i);
	for (int i = 0, q = n - 1; i < n; i++) {  // 偶数：中心对 i<j，q 单调左移
		while (q > i && a[i] + a[q] > Y) q--;
		if (q <= i) break;
		res += CC(n - 1, i + 1) - CC(i + n - q - 1, i + 1) + mod; // j in (i, q] 朱世杰合并
	}
	return res % mod;
}
void Solve() {
	int n; ll x; cin >> n >> x;
	vector<ll> a(n); for (ll &v : a) cin >> v;
	sort(all(a));
	cout << (F(a, 2 * x) - F(a, 2 * x - 1) + mod) % mod << '\n';  // 恰好 = 不超过 2x 减不超过 2x-1
}
// 对拍：与官方 std、Python 暴力随机 3000 组一致
		\end{minted}

		\subsection{常用组合恒等式}

		\begin{itemize}
			\item 对称：$\binom{n}{k} = \binom{n}{n-k}$；
			\item 吸收：$k\binom{n}{k} = n\binom{n-1}{k-1}$
			      ——见到组合数前面乘着下指标就用它；
			\item 二项式定理的两个特例：
			      $\sum_{i}\binom{n}{i} = 2^n$，
			      $\sum_{i}(-1)^i\binom{n}{i} = [n = 0]$；
			\item 带权和（都由吸收恒等式导出）：
			      \[
			        \sum_{i}i\binom{n}{i} = n2^{n-1},\quad
			        \sum_{i}i^{2}\binom{n}{i} = n(n+1)2^{n-2};
			      \]
			\item 上指标求和即朱世杰恒等式（见上文）；
			\item 平方和即范德蒙德卷积：$\sum_i\binom{n}{i}^2 = \binom{2n}{n}$
			      （「左右各选 $r$ 个」与平移版见范德蒙德卷积一节）；
			\item 奇偶项和：$\sum_i\binom{n}{2i} = \sum_i\binom{n}{2i+1} = 2^{n-1}$（$n\ge1$）；
			\item 行前缀和 $S(n,k)=\sum_{i\le k}\binom{n}{i}$ 没有封闭式，但能 $O(1)$ 移动：
			      \[
			        \begin{aligned}
			        	S(n+1,k) &= 2S(n,k)-\binom{n}{k},\\
			        	S(n,k+1) &= S(n,k)+\binom{n}{k+1},
			        \end{aligned}
			      \]
			      多组询问 $(n,k)$ 用莫队；
			\item 交错前缀和：$\sum_{i\le k}(-1)^i\binom{n}{i} = (-1)^k\binom{n-1}{k}$；
			\item 三项式换位：$\binom{n}{k}\binom{k}{m} = \binom{n}{m}\binom{n-m}{k-m}$
			      ——「先选 $k$ 个再从中挑 $m$ 个」交换求和顺序时用；
			\item 上指标卷积：$\sum_{k}\binom{k}{a}\binom{n-k}{b} = \binom{n+1}{a+b+1}$
			      （枚举第 $a+1$ 个被选元素的位置）；
			\item 斜线和：$\sum_i\binom{n-i}{i} = F_{n+1}$（$F_1=F_2=1$，即 $1\times2$ 骨牌铺 $1\times n$）；
			\item 多重集排列 $\dfrac{n!}{\prod a_i!}$，环排列 $(n-1)!$（项链可翻转再 $\div 2$，$n\ge3$）；
			\item 有上界的插板：$x_1+\cdots+x_k=n$，$0\le x_i\le c$ 的解数，容斥钦定 $j$ 个超界：
			      \[
			        \sum_{j:\,j(c+1)\le n}(-1)^j\binom{k}{j}\binom{n-j(c+1)+k-1}{k-1};
			      \]
			\item 格路反射：$(0,0)\to(n,m)$ 只走右 / 上、\textbf{不碰}直线 $y=x+t+1$（$m\le n+t$）的路径数
			      $\binom{n+m}{n}-\binom{n+m}{n+t+1}$（起点关于该直线反射到 $(-t-1,t+1)$）；$t=0,\,m=n$ 即 Catalan；
			\item 错排递推：$D_n=(n-1)(D_{n-1}+D_{n-2})$，$D_1=0,\ D_2=1$。
		\end{itemize}



\end{document}
